problem:  
Let \( (M^{4n+3}, g, \mathcal{S}) \) be a **regular**, compact, simply connected 3‑Sasakian manifold with structure \(\mathcal{S} = \{(\xi_i, \eta_i, \Phi_i)\}_{i=1}^3\)\, generating an effective isometric \(\mathrm{Sp}(1)\)-action by its Reeb fields.  
Let \(\pi: M \to B\) denote the projection onto the quaternionic‑Kähler base \(B = M/\mathrm{Sp}(1)\), which is smooth by regularity.  

It follows from results of Boyer–Galicki and Swann that for regular 3‑Sasakian manifolds, every isometry of the base \(B\) lifts (up to the \(\mathrm{Sp}(1)\) action) to a 3‑Sasakian automorphism of \(M\): hence there is a canonical embedding  
\[
\mathrm{Sp}(1)\times\mathrm{Iso}(B) \hookrightarrow \mathrm{Aut}(M,\mathcal{S}) \subseteq \mathrm{Iso}(M,g).
\]  
All known examples where \(\mathrm{Aut}(M,\mathcal{S})\) strictly exceeds this canonical group correspond to *homogeneous* 3‑Sasakian manifolds arising from Wolf spaces.  

**Open problem:**  
Does there exist a compact, simply connected, **regular non‑homogeneous** 3‑Sasakian manifold \( (M^{4n+3}, g, \mathcal{S}) \) for which  
\[
\mathrm{Aut}(M,\mathcal{S})
\text{ strictly contains }
\mathrm{Sp}(1) \times \mathrm{Iso}(B),
\quad \text{but is strictly smaller than } \mathrm{Iso}(M,g)?
\]  
Equivalently, can one construct (or rule out) a regular 3‑Sasakian manifold with an **intermediate automorphism group**—larger than that induced by the quaternionic‑Kähler base, yet still insufficient to make \(M\) homogeneous—thereby breaking the known “two‑level” pattern of symmetry in current examples?

Why is it a "good" problem:  
This problem isolates a fundamental question about **symmetry rigidity** in 3‑Sasakian geometry. It draws upon the precise Boyer–Galicki framework, where known examples admit only two extremes of symmetry: (i) maximally symmetric homogeneous models linked to Wolf spaces, and (ii) deformations with sharply reduced symmetry. Whether intermediate symmetry tiers exist is currently unestablished in the literature—even for regular cases—with no known construction yielding them and no rigidity theorem ruling them out.  

Resolving this would clarify how **quaternionic‑Kähler base symmetries lift through the 3‑Sasakian fibration**, potentially identifying new geometric mechanisms beyond the homogeneous paradigm. It also connects differential geometry, Lie group actions, and special holonomy theory in a single, concise question: a “narrow entrance, profound depth” problem whose answer would reshape understanding of the relationship between symmetry and special Riemannian structures.